\section{External Validity}
\label{sec:external}

\subsection{Evaluation of Public Corpora}

We surveyed candidate public cybersecurity datasets to evaluate the feasibility of
instantiating the five-node ZTA graph on real-world enterprise telemetry, yielding two
key findings:

\textbf{Lack of Zero Trust benchmarks.} Existing Zero Trust literature on continuous
trust scoring and policy decision points evaluates predominantly on proprietary or
unreleased simulated environments, with no standardized benchmark for continuous authorization.

\textbf{Absence of concurrent identity and TLS fingerprints.} Authentication graph
corpora (such as LANL~\cite{lanl}, DARPA OpTC~\cite{optc}, and CERT~\cite{cert}) provide
extensive event volumes and lateral movement labels, but omit client TLS fingerprints entirely.
Conversely, network testbeds providing packet captures (such as AIT~\cite{ait} and
CIC-IDS2017~\cite{cicids}) enable TLS fingerprint extraction, but either lack user identity
associations (precluding principal-centric modeling) or omit labeled lateral movement activity.

One public dataset provides the required intersection of telemetry fields.

\subsection{PicoDomain}

PicoDomain~\cite{picodomain} is a \PicoSpan-day capture of an enterprise Windows domain
intrusion across five workstations and a domain controller, published as Zeek JSON logs
with an annotated red-team activity log. It provides real telemetry covering all five node
types in our graph schema: its \texttt{ssl.log} contains \PicoSSLrecords{} records with
\texttt{ja3} fingerprints over \PicoJAthree{} distinct client configurations, while
\texttt{kerberos.log} and \texttt{ntlm.log} capture both human user principals and machine accounts.

\begin{table}[t]
\centering
\caption{PicoDomain mapped onto the five-node chain schema.}
\label{tab:picomap}
\footnotesize
\begin{tabular}{@{}ll@{}}
\toprule
Node & Zeek field \\
\midrule
source   & \texttt{id.orig\_h} \\
config   & \texttt{ssl.log:ja3} \\
device   & Kerberos machine account \texttt{HOST\$}, \texttt{ntlm:hostname} \\
user     & \texttt{kerberos:client} (non-machine), \texttt{ntlm:username} \\
resource & \texttt{smb\_mapping:path}, \texttt{smb\_files:name}, \\
         & \texttt{http:uri}, \texttt{dce\_rpc:endpoint} \\
\bottomrule
\end{tabular}
\end{table}

The red-team exercise covers multi-stage lateral progression: initial compromise via a
trojanized archive, privilege escalation, credential dumping via \emph{Mimikatz}, domain
reconnaissance, remote execution across three additional hosts using DCOM and WMI, and
interactive command execution under an alternate domain account. Both lateral movement
and credential reuse are present in the activity trace.

\subsection{Telemetry Reconstruction and Limitations}

In network monitoring telemetry, TLS handshakes and authentication transactions typically
occur over separate transport connections. In PicoDomain's Zeek logs, connection identifiers
(\texttt{uid}) between \texttt{ssl} and authentication protocols (\texttt{kerberos}/\texttt{ntlm})
exhibit zero overlap (\PicoUidOverlap). Consequently, the \textit{config}\,$\rightarrow$\,\textit{user}
binding is reconstructed temporally: for each client source address, the most recently observed
authenticated principal and TLS fingerprint within an active session window are associated.

The temporal binding window is set to \PicoTTL~s, corresponding to the default Kerberos ticket
granting service lifetime. Binding coverage across entities is sensitive to this parameter:
user binding coverage achieves \PicoCovUserShortest{} at 900\,s, \PicoCovUserShort{} at 3600\,s,
and \PicoCovUser{} at \PicoTTL~s, with device and configuration coverage reaching \PicoCovDevice{}
and \PicoCovConfig{} respectively. Unattributed events map to dedicated per-address sentinel nodes
rather than a shared global identifier, preventing memory state contamination across distinct endpoints.

The resulting stream comprises \PicoEvents{} access events over \PicoNodes{} unique nodes
(\PicoUsers{} users, \PicoDevices{} devices, \PicoSources{} sources, \PicoConfigs{} configurations,
and \PicoResources{} resources), with \PicoNlateral{} lateral movement, \PicoNtheft{} credential
reuse, and \PicoNctx{} contextual anomaly events (\PicoAnomFrac{} anomaly ratio). Ground-truth labels
are assigned by matching access events against red-team activity logs within a $\pm90$\,s window
on host or principal identity. The red-team activity log timestamps carry an estimated uncertainty
of $\pm1$ minute, reflecting manual operator logging.

With eight distinct source addresses over \PicoSpan{} days, PicoDomain serves as an empirical
case study rather than an enterprise-scale benchmark. Variance across the \PicoNlateral{} lateral
events is inherently non-trivial, and policy violations cannot be evaluated due to the absence
of authorization policy definitions. Nonetheless, it provides valuable empirical confirmation
that the five-node graph schema maps directly to enterprise monitoring telemetry.

\subsection{Model results on PicoDomain}
\label{sec:pico-model}

We evaluated the one-class temporal learning pipeline (\texttt{tests/eval\_picodomain.py})
on the mapped PicoDomain stream (\PicoModelSettings) without per-corpus hyperparameter tuning.
Ranking performance demonstrates discrimination across all three evaluable attack classes:
aggregate AUC \PicoModelAggAuc{} (AP \PicoModelAggAp), lateral movement AUC \PicoModelLatAuc{}
(AP \PicoModelLatAp, $n=\PicoNlateral$), credential reuse AUC \PicoModelTheftAuc{} (AP
\PicoModelTheftAp, $n=\PicoNtheft$), and contextual anomaly AUC \PicoModelCtxAuc{} (AP
\PicoModelCtxAp, $n=\PicoModelCtxN$). While absolute performance is lower than on the synthetic
benchmark, reflecting the limited sample size and small entity population, the model successfully
maintains inductive ranking ability on real network traces.

We report ranking metrics (AUC and AP) rather than operating recall for this dataset due to the
temporal distribution of attacks. In PicoDomain, red-team activities occur exclusively in the
final segment of the capture (commencing after 84\% of the event stream). Under a strict
chronological 70/10/20 split, the validation slice contains zero lateral movement or credential
theft instances. Threshold calibration for specific recall operating points requires positive
validation instances; calculating recall without validation instances would introduce threshold
bias unrepresentative of actual detector capability.

\subsection{LANL as an Architectural Ablation}

While LANL~\cite{lanl} is an established benchmark for lateral movement on authentication graphs,
it lacks client application fingerprints, device attestation, and resource sensitivity metadata.
Evaluating our five-node architecture on LANL necessarily requires ablating the client-configuration
and device nodes, collapsing the schema onto a bipartite computer-to-computer graph. While evaluating
this bipartite degradation provides a useful structural ablation, comparing LANL numbers directly
against systems designed exclusively for host authentication logs would evaluate the architecture
with its core contextual components disabled.

\subsection{Future Corpus Directions}

Future empirical validation on enterprise-scale data can build upon corpora such as the AIT
Log Data Set~\cite{ait}, which features multiple testbeds with real network boundaries, diverse
authentication mechanisms, and packet captures from which TLS fingerprints can be derived.
Providing comprehensive lateral movement ground truth on such multi-source telemetry represents
an important direction for the Zero Trust research community.
